\documentclass[aspectratio=169]{beamer}

\usetheme{Madrid}
\usecolortheme{beaver}
\setbeamertemplate{navigation symbols}{}
\setbeamertemplate{footline}[frame number]

\usepackage[utf8]{inputenc}
\usepackage{booktabs}
\usepackage{array}
\usepackage{tabularx}

\title[Hoot: Progress Report 1]{Hoot: Progress Report 1}
\subtitle{Predictive Household Inventory Tracker}
\author{Michael Hannan}
\date{Week 9 --- October 22, 2026}

\begin{document}

\begin{frame}
  \titlepage
\end{frame}

% ---------------------------------------------------------
\begin{frame}{Where We Are}
  \begin{itemize}
    \item This is the Week~9 progress check from the project breakdown.
    \item Original plan for today: basic database, item/purchase save-and-load, and a simple dashboard.
    \item What works: a working app demo in the browser first.
    \item The server and database are started, but not finished yet.
    \item This talk covers three things: \textbf{Done}, \textbf{TODO}, and \textbf{Update}.
  \end{itemize}
\end{frame}

% ---------------------------------------------------------
\begin{frame}{Task Done --- What You Can Use Today}
  \begin{itemize}
    \item \textbf{Dashboard} --- items sorted by how soon they may run out, with clear status labels.
    \item \textbf{Restock calendar} --- a two-week view of upcoming restock dates.
    \item \textbf{Add item / Log purchase} --- simple forms to track staples and record what you bought.
    \item \textbf{Item page} --- purchase history, edit or delete, and a chart of the predicted restock.
    \item \textbf{One-time list} --- a checklist for one-off needs that should not be predicted.
  \end{itemize}
\end{frame}

% ---------------------------------------------------------
\begin{frame}{Task Done --- How Predictions Work (For Now)}
  \begin{itemize}
    \item New items use a manual guess: ``I usually buy this every X days.''
    \item After three purchases, Hoot learns the gap between buys and estimates the next date.
    \item Demo data (milk, eggs, paper towels, and more) is saved in the browser.
    \item The UI is set up so it can later talk to a real server without a full redesign.
    \item Planning and design docs are written: breakdown, project plan, product notes, and UI flows.
  \end{itemize}
\end{frame}

% ---------------------------------------------------------
\begin{frame}{Task Done --- By Project Area}
  \small
  \begin{tabularx}{\linewidth}{>{\raggedright\arraybackslash}p{3.1cm} X}
    \toprule
    \textbf{Area} & \textbf{Status} \\
    \midrule
    Frontend &
      Mostly done: dashboard, forms, calendar, item page, one-time list \\[4pt]
    Prediction &
      Works in the browser demo; not yet saved on the server \\[4pt]
    Backend / API &
      Project started; real item and purchase endpoints still needed \\[4pt]
    Database &
      Tables are planned; not built in PostgreSQL yet \\[4pt]
    Barcode lookup &
      Placeholder on the item page; no live product lookup yet \\[4pt]
    Testing \& Deploy &
      Manual demo testing only; no packaging or deploy setup yet \\
    \bottomrule
  \end{tabularx}
\end{frame}

% ---------------------------------------------------------
\begin{frame}{Task TODO --- Next (Weeks 10--12)}
  \begin{itemize}
    \item Build the database tables for users, items, purchases, and predictions.
    \item Finish the server API so items and purchases can be saved for real.
    \item Connect the front end to that API instead of browser-only storage.
    \item Move prediction math onto the server and keep a history of each estimate.
    \item Add barcode lookup so scanning can fill in product details.
  \end{itemize}
\end{frame}

% ---------------------------------------------------------
\begin{frame}{Task TODO --- Later}
  \begin{itemize}
    \item Login and register screens.
    \item Check how accurate predictions were over time.
    \item More testing, bug fixes, and optional Docker deploy.
    \item Nice-to-haves later: shared households, shopping mode, price trends, reminders.
    \item Final week: demo slides, demo video, final report, and code hand-in.
  \end{itemize}
\end{frame}

% ---------------------------------------------------------
\begin{frame}{Task Update --- Plan vs.\ Reality}
  \begin{itemize}
    \item \textbf{Ahead on the UI:} we planned a basic dashboard; we have ranking, calendar, charts, and full add/edit/delete flows.
    \item \textbf{Behind on the database:} save-and-load on the server was due now; that work comes next.
    \item \textbf{Prediction early, but local:} the learning math works in the demo; it still needs to live on the server.
    \item \textbf{One-time list early:} that checklist was planned later and is already in the demo.
  \end{itemize}
\end{frame}

% ---------------------------------------------------------
\begin{frame}{Task Update --- What Comes Next}
  \begin{itemize}
    \item First job: finish the database and server so the app does not rely on browser storage.
    \item Next: barcode and product lookup once the core tables exist.
    \item Goal for Progress Report~2 (Nov~10): live API, server predictions, and barcode support started.
    \item Watch-out: keep the same prediction rules when moving from demo to server.
  \end{itemize}
\end{frame}


\end{document}
